
% Arabic abstract (bilingual thesis requirement).
%
% This file is a standalone XeLaTeX document. The main thesis document
% (main.tex) is built with pdflatex + mathptmx (packages.tex) and does not
% support Arabic script, so this abstract is compiled separately with:
%
%     xelatex abstract_ar.tex
%
% and the resulting single-page abstract_ar.pdf is renamed to
% arabicabstract.pdf and included in the main document via
% \includepdf in titlepage.tex.
%
% Required packages: fontspec, polyglossia, and the Amiri Arabic font
% (Ubuntu/Debian: texlive-lang-arabic, fonts-hosny-amiri).

\documentclass[12pt,a4paper,oneside]{report}

\usepackage{geometry}
\newgeometry{
    left=2.5cm,
    right=3.0cm,
    top=2.0cm,
    bottom=2.5cm
}

\usepackage{fontspec}
\usepackage{polyglossia}

\setdefaultlanguage{arabic}
\setotherlanguage{english}

\newfontfamily\arabicfont[Script=Arabic]{Amiri}

\usepackage{setspace}
\doublespacing

\pagestyle{empty}

\begin{document}

% ============================================================
% Thesis title
% ============================================================

\begin{center}

{\Large\bfseries
FusionRAG: إطار حوار متعدد الوسائط وواعٍ بالسياق، يدمج التعرف الآلي على الكلام
(ASR) والإشارات البصرية والنماذج اللغوية الكبيرة المعزَّزة بالاسترجاع
}

\vspace{0.3cm}

{\large
وفاء عبدالله العمري
}

\vspace{0.2cm}

{\normalsize
المشرف: أ.د. قيصر عباس
}

%\vspace{0.8cm}

%{\Large\bfseries
%المستخلص
%}

\end{center}

\vspace{0.1cm}

% ============================================================
% Arabic Abstract
% ============================================================

\noindent
تتطلب أنظمة الحوار المنطوق القائمة على نماذج اللغة الكبيرة تحقيق متطلبين غالبًا ما يُدرَسان
بمعزل عن بعضهما: الإدراك الموثوق للمدخلات الصوتية في ظل ظروف صوتية وبصرية واقعية، وارتكاز
الاستجابات المولَّدة على أدلة يمكن التحقق منها بدلًا من الاعتماد على المعرفة المخزَّنة داخل
النموذج أثناء التدريب المسبق. تتناول هذه الأطروحة مشكلة الدمج بين هذين المتطلبين من خلال
اقتراح إطار عمل حواري نمطي يُعرف بـ FusionRAG، يربط بين نظام فرعي للإدراك الكلامي متعدد
الوسائط (MSP) ونظام فرعي للاستدلال المعزَّز بالاسترجاع عبر واجهة واحدة تتمثل في النص
المُتعرَّف عليه، بما يتيح تطوير كل نظام فرعي وتدريبه وتقييمه بشكل مستقل وفق معايير مرجعية
مناسبة لمهمته. يتعرّف النظام الفرعي للإدراك الكلامي على المحتوى المنطوق عبر ثلاثة مسارات:
صوتي فقط، وبصري فقط، وسمعي-بصري مدمج، حيث يُبنى المسار الأخير باستخدام آلية دمج قائمة على
انتباه متقاطع ثنائي الاتجاه تُدرَّب بشكل مشترك مع المسارين الأحاديي الوسيط تحت هدف تدريب
موحّد قائم على تصنيف زمني اتصالي (CTC). وبمقارنته بنظامين سمعيين-بصريين قويين وسابقي النشر،
هما Auto-AVSR و Whisper-Flamingo، أظهر نموذج الدمج المقترح تفوقًا واضحًا في متانة الأداء
تحت التدهور الصوتي: فعلى مجموعة اختبار LRS2، حقق أدنى معدل خطأ في الكلمات بين الأنظمة
الثلاثة تحت أشد ظروف التشويش متعدد المتحدثين، رغم تدريبه على بيانات سمعية-بصرية أقل بكثير؛
وعلى معيار AVCocktail الواقعي متعدد المتحدثين، حقق أدنى معدل خطأ تحت كل من تجزئة الكشف عن
المتحدث النشط والتجزئة الثابتة، مع أقل تدهور إجمالي بين استراتيجيات التجزئة من بين الأنظمة
الثلاثة، مما يؤكد أن الدمج متعدد الوسائط ثنائي الاتجاه يحقق مكسبًا حقيقيًا في متانة
الإدراك وليس مجرد مجاراة للأنظمة السابقة تحت الظروف المواتية فقط. أما النظام الفرعي
للاستدلال المعزَّز بالاسترجاع، المبني بالكامل من مكوّنات استرجاع وإعادة ترتيب وتوليد
مُدرَّبة سلفًا، فيستقبل النص المُتعرَّف عليه كاستعلام ويولّد استجابة مقيَّدة بالأدلة
المسترجَعة؛ وقد حقق استدعاء استرجاع تجاوز أربعة وتسعين بالمئة عند عمق عشرة مقاطع عبر ثلاثة
مجالات متباينة من معيار RAGBench، فيما حسّنت إعادة الترتيب باستخدام المُرمِّز المتقاطع جودة
الاسترجاع والتوليد في المجالات المتوافقة مع بيانات تدريبه، وكشفت عن حد دقيق يمكن إرجاعه
إلى اختلاف المجال في مجال طبي حيوي ذي مفردات مختلفة بشكل ملحوظ. وتُثبت هذه النتائج مجتمعةً
أن بنية الدمج المقترحة تحقق تحسينات ملموسة وقابلة للقياس في متانة الإدراك مقارنة بالأنظمة
السمعية-البصرية المرجعية تحت ظروف محادثة صعبة، وأن نظام الاستدلال المعزَّز بالاسترجاع
المقترح يولّد استجابات يمكن التحقق من ارتكازها على الأدلة المسترجَعة لها، بما يُثبت أن
FusionRAG يمثّل بنية نمطية فعّالة ومُقيَّمة بدقة لأنظمة الحوار المنطوق المرتكزة على
المعرفة.

\end{document}
